% ============================================================
% UGA-Themed Beamer Template (Clean MWE)
% ============================================================

\documentclass[aspectratio=169,professionalfonts]{beamer}

% ---------- COLORS (swap with official UGA hex if you like) ----------
\usepackage{xcolor}
\definecolor{UGARed}{HTML}{BA0C2F}   % Verify with brand guide
\definecolor{UGABlack}{HTML}{000000}
\definecolor{UGAGray}{HTML}{8A8D8F}

% ---------- THEME & PROGRESS DOTS ----------
\usetheme{default}
\useinnertheme{circles}
\useoutertheme[subsection=false]{miniframes}
\usefonttheme{professionalfonts}
\setbeamertemplate{navigation symbols}{}
% Head/foot colors
\setbeamercolor{subsection in head/foot}{fg=white,bg=UGARed}
\setbeamercolor{headline}{fg=white, bg=UGARed}

\setbeamercolor{section in head/foot}{fg=white,bg=UGARed}
\setbeamercolor{section in head/foot shaded}{fg=white,bg=UGARed}
\setbeamercolor{footline}{bg=UGABlack,fg=white}
\setbeamercolor{frametitle}{fg=UGABlack}
\setbeamercolor{title}{fg=UGABlack}
\setbeamercolor{normal text}{fg=UGABlack}
\setbeamercolor{structure}{fg=UGARed}

% --- HEADLINE: slide dots grouped by section ---
\setbeamersize{mini frame size=0.25cm,mini frame offset=0pt}
% Show section labels above the clickable slide dots.
\setbeamertemplate{section in head/foot}[default]
\setbeamertemplate{section in head/foot shaded}[default][75]
\setbeamertemplate{mini frame}{%
  \tikz[baseline=-0.5ex] \filldraw[draw=white,fill=white,line width=0.6pt] (0,0) circle (2pt);%
}
\setbeamertemplate{mini frame in current subsection}{%
  \tikz[baseline=-0.5ex] \filldraw[draw=white,fill=UGARed,line width=0.6pt] (0,0) circle (2pt);%
}
\setbeamertemplate{mini frame in other subsection}{%
  \usebeamertemplate{mini frame in current subsection}%
}
\setbeamertemplate{mini frame in other section}{%
  \usebeamertemplate{mini frame in current subsection}%
}
\setbeamertemplate{headline}{%
  \begin{beamercolorbox}[wd=\paperwidth]{headline}
    \vskip1ex
    \insertnavigation{\paperwidth}
    \vskip3ex
  \end{beamercolorbox}
}

% Simple footline with slide numbers
\setbeamertemplate{footline}{
  \leavevmode%
  \hbox{%
    \begin{beamercolorbox}[wd=\paperwidth,ht=2.5ex,dp=1ex,leftskip=1em,rightskip=1em]{footline}%
      \insertshortauthor{} \hfill \insertframenumber/\inserttotalframenumber
    \end{beamercolorbox}}%
}

% ---------- PACKAGES ----------
\usepackage{amsmath,amssymb,mathtools}
\usepackage{graphicx}
\usepackage{booktabs}
\usepackage{bbm}
\usepackage{array}
\usepackage{adjustbox}
\usepackage{hyperref}
\usepackage{tikz}
\usetikzlibrary{arrows.meta, positioning}
\usepackage{tcolorbox}
\usepackage{eurosym}
\usepackage{threeparttable}
% ---------- METADATA (fill these later) ----------
\title{Optimal Dynamic Pricing With Preferences for Novelty}
\subtitle{}
\author{Tate Mason}
\institute{The University of Georgia}
\date{\today}

% ---------- MACROS ----------
% Numbered equation with label
\newcommand{\beq}[1]{\begin{equation}\label{#1}}
\newcommand{\eeq}{\end{equation}}
% Accent color highlight
\newcommand{\hl}[1]{\textbf{\textcolor{UGARed}{#1}}}
% Simple figure placeholder macro
\newcommand{\figph}[2][0.9\linewidth]{%
  \begin{figure}
    \centering
    \fbox{\rule{0pt}{0.55#1}\rule{#1}{0pt}}%
    \caption{#2}
  \end{figure}
}

% ============================================================

% Keep navigation links in the theme's white foreground on the red header.
\hypersetup{colorlinks=true, linkcolor=., urlcolor=UGARed}

% ---------- EXERCISE AND ESTIMATION FILE PATHS ----------
\newcommand{\supplyexercisepath}{../pricing_exercises/run_supply_exercises.py}
\newcommand{\pricingexercisepath}{../pricing_exercises/run_empirical_targeting.py}
\newcommand{\supplyresultspath}{../pricing_exercises/supply_outputs/}
\newcommand{\pricingresultspath}{../../../shell_files/pricing_ex/}
\newcommand{\estimationnocfpath}{../../../shell_files/res_het_nocf/}
% No res_het directory is present; the available heterogeneous fit is res_het_hist.
\newcommand{\estimationhetpath}{../../../shell_files/res_het_hist/}
\newcommand{\gammadistributionnocfpath}{../../../shell_files/res_het_nocf/history_distributions.pdf}
\newcommand{\gammadistributionhetpath}{../../../shell_files/res_het_hist/history_distributions.pdf}

\begin{document}

\begin{frame}[plain,noframenumbering]
  \titlepage
\end{frame}

\section{Introduction}
% Each frame title is independent of the section name in the navigation bar.
\begin{frame}{Motivation}
  la
\end{frame}
\begin{frame}{Research Question}
  la
\end{frame}

\section{Motivation from Data}

\section{Supply Theory}
\begin{frame}{Supply Setting}
  \begin{itemize}
    \item Monopolistic firm producing a single product $j$
    \item Firm is profit maximizing, setting prices according to that objective
    \item A consumer is born into the world with a type $r\in\{H,L\}$, determining how much utility they derive from switching and an initial state $h_0\in\{0,j\}$
    \item Consumers live for two periods, purchasing the inside good or opting for the outside good in each period
    \item The firm knows initial states and types, setting prices to individual consumers by type and history
  \end{itemize}
\end{frame}

\begin{frame}{Consumer's Problem}
  \begin{tcolorbox}[
    boxsep=0pt,
    left=2pt, right=2pt,
    top=1pt, bottom=1pt,
    before skip=2pt, after skip=2pt
    ]
    \setlength{\abovedisplayskip}{2pt}
    \setlength{\belowdisplayskip}{2pt}
    \begin{gather*}
      V^{inside}_{jt} = \underbrace{\beta}_{\text{Constant}} + \overbrace{\gamma_r\mathbbm{1}\{h=0\}}^{\text{Utilty from Switching}} - \underbrace{\alpha p_t(h)}_{\text{Price Disutility}}, \ V^{Outside}_{0t} = \underbrace{\gamma_r\mathbbm{1}\{h=j\}}_{\text{Utility from Switching}} \\
      v_t(x) = V_{jt}^{inside} - V_{0t}^{outside} = \beta + \underbrace{\gamma[\mathbbm{1}\{h=0\} - \mathbbm{1}\{h=j\}]}_{q(h)} - \alpha p_t(h)
    \end{gather*}
  \end{tcolorbox}
  \begin{itemize}\setlength{\itemsep}{2pt}
    \item Prices, set by the firm, and switching utility depend on history $h=x_{t-1}$
    \item Utility for the outside good is normalized to zero
  \end{itemize}
\end{frame}

\begin{frame}{Firm's Problem}
  \begin{tcolorbox}
    [
    boxsep=0pt,
    left=2pt, right=2pt,
    top=1pt, bottom=1pt,
    before skip=2pt, after skip=2pt
    ]
    \setlength{\abovedisplayskip}{2pt}
    \setlength{\belowdisplayskip}{2pt}
    \begin{gather*}
      \pi_t(h) = \max_{p_t(h)} s_t(h)[p_t(h) - c] \\ 
      \Pi(h) = \underbrace{s_1(h)(p_1(h) - c)}_{\text{Margin in Period 1}} + \underbrace{\delta\{s_1(h)\pi_2(j) + (1-s_1(h))\pi_2(0)\}}_\text{{Continuation Value}} 
    \end{gather*}
  \end{tcolorbox}
  \begin{itemize}
    \item Firm profits are comprised of period one margins, given history of consumption, and the value of that consumer in the different states in period two
    \item The firm will set prices according to consumer type and history to maximize expected profit
  \end{itemize}
\end{frame}

\begin{frame}{Optimal Terminal Price and Profit}
  \begin{tcolorbox}
    [
    boxsep=0pt,
    left=2pt, right=2pt,
    top=1pt, bottom=1pt,
    before skip=2pt, after skip=2pt
    ]
    \setlength{\abovedisplayskip}{2pt}
    \setlength{\belowdisplayskip}{2pt}
    \begin{align*}
      \frac{\partial\pi_2(h_1)}{\partial p_2(h_1)} &= s_2(h_1) + (p_2 - c)\frac{\partial s_2(h_1)}{\partial p_2(h_1)} \\
      p_2^*(h_1)&= c + \frac{1}{\alpha(1-s_2^*(h_1))} \\
      \pi^*_2(p_2^*(h_1); h_1) &= \frac{s_2^*(h_1)}{\alpha(1-s_2^*(h_1))}
    \end{align*}
  \end{tcolorbox}
  \begin{itemize}
    \item Optimal period two price is cost minus a demand markup term
    \item Period two profit at optimal prices is the odds of purchase, scaled by the inverse price sensitivity
  \end{itemize}
\end{frame}

\begin{frame}{First Period Optimization}
  \begin{tcolorbox}
    [
    boxsep=0pt,
    left=2pt, right=2pt,
    top=1pt, bottom=1pt,
    before skip=2pt, after skip=2pt
    ]
    \setlength{\abovedisplayskip}{2pt}
    \setlength{\belowdisplayskip}{2pt}
    \begin{align*}
      \frac{d\Pi}{d p_1(h_0)} &= s_1 + \frac{\partial}{\partial p_1(h)}s_1(h_0)[p_1(h_0) - c + \delta(\pi_2(j) - \pi_2(0))] \\
      p_1^*(h_0) &= c + \frac{1}{\alpha(1-s_1^*(h_0))} - \delta(\pi_2(j) - \pi_2(0)) \\
      \Pi^*(h) &= \delta\pi_2(0) + \frac{s_1^*(h_0)}{\alpha(1-s_1^*(h_0))}
    \end{align*}
  \end{tcolorbox}
\end{frame}

\begin{frame}{Optimal Profit Given Switching Type}
  \begin{tcolorbox}
    [
    boxsep=0pt,
    left=2pt, right=2pt,
    top=1pt, bottom=1pt,
    before skip=2pt, after skip=2pt
    ]
    \setlength{\abovedisplayskip}{2pt}
    \setlength{\belowdisplayskip}{2pt}
    \begin{align*}
      \frac{d\Pi}{d\gamma} &= \frac{q(h)s_1^*(h) - \delta s_1^*(h)s_2^*(j) + \delta(1-s_1^*(h))s_2^*(0)}{\alpha}
    \end{align*}
  \end{tcolorbox}
  \begin{itemize}
    \item There are two interpretations:
    \begin{enumerate}
      \item If $h_0=j$, $q(h)<0$, leading us to ambiguous signage as we raise profits from initially outside consumers but see opposing forces for the initially inside: leave in period one, return in period two 
      \item If $h_0=0$, $q(h)>0\Rightarrow \frac{d\Pi}{d\gamma} < 0$, meaning that profits decrease as $\gamma$ increases
    \end{enumerate}
    \item Points to firms potentially reaping great benefit from knowing consumer types and histories
  \end{itemize}
\end{frame}

\begin{frame}{Pricing Simulations Specification}
  \centering
  \small
  \begin{tabular}{lll}
    \toprule
    Parameter & Description & Value \\
    \midrule
    $\beta$ & Inside-good intercept & $2$ \\
    $\alpha$ & Price sensitivity & $1$ \\
    $\delta$ & Discount factor & $0.92$ \\
    $c$ & Marginal cost & $0$ \\
    $\gamma$ & Switching preference & $[-3,3]$ \\
    \midrule
    \multicolumn{2}{l}{Initial Consumption State}
      & $h_0\in\{0,j\}$ \\
    \multicolumn{2}{l}{Feasible Prices}
      & $p_t(h)\in\mathbb{R}$ \\
    \bottomrule
  \end{tabular}
  \begin{itemize}
    \item Allow for negative pricing
    \item Firm has knowledge of both type and history, allowing for targeted pricing
    \item Four consumption histories possible $(h_0, h_1) \in \{(0,0), (0,j), (j,j), (j,0)\}$
  \end{itemize}
\end{frame}

\section{Appendix}
  \begin{frame}{Shares and Elasticities}
  \begin{tcolorbox}
    [
    boxsep=0pt,
    left=2pt, right=2pt,
    top=1pt, bottom=1pt,
    before skip=2pt, after skip=2pt
    ]
    \setlength{\abovedisplayskip}{2pt}
    \setlength{\belowdisplayskip}{2pt}
    \begin{gather*}
      s_t(h) = \frac{e^{v_t(h)}}{1+e^{v_t(h)}}, \ s_{0t}(h) = 1-s_t(h) \\
      \frac{\partial s_t(h)}{\partial p_t(h)} = -\alpha s_t(h)(1-s_t(h)) \\
      \frac{\partial s_t(h)}{\partial\gamma} = q(h)s_t(h)(1-s_t(h))
    \end{gather*}
  \end{tcolorbox}
  \end{frame}

\end{document}
